%% kapa: exam=1 round=27 year=2016 period=1 form=A subject=real_estate subject_ko=부동산학원론 seq=10 number=90
%% source: Q-Net 감정평가사 기출문제 (cst00309) artlSeq=5169402 file=extracted/1008/5169402/제27회 감정평가사 1차 문제/제27회 감정평가사 1교시문제(A형).pdf
%% conversion: kordoc 4.24.51 -> tools/gen1.py (auto); review=proofread-sonnet; flags=none
\begin{question}{90}
디파스퀠리\&위튼(DiPasquale\&Wheaton)의 4사분면 모형에 관한 설명으로 옳지 않은 것은?
\begin{choices}
  \item 부동산 공간시장과 부동산 자산시장의 관계를 설명한 모형이다.
  \item 1사분면은 부동산 가격과 공간재고량의 관계를 나타낸다.
  \item 2사분면은 부동산 가격과 임대료의 관계를 나타낸다.
  \item 3사분면은 부동산 가격과 신규건설량의 관계를 나타낸다.
  \item 4사분면은 신규건설량과 공간재고량의 관계를 나타낸다.
\end{choices}
\end{question}
\answer{2}
